% IEEE Access — rebuilt manuscript (Access-2026-27906).
% DRAFT 2026-07-20. Result-independent sections written now; every grid number is a VISIBLE
% placeholder \gridnum{...} that renders as [[GRID: ...]] so it can NEVER be mistaken for a final
% value (DECISIONS D35: all result numbers trace to the 740-patient grid, never the 250 sanity
% stage). Do NOT fill Results or the calibration-mechanism numbers until the grid is complete and
% verify_grid_consistency passes.
\documentclass{article}
\usepackage{amsmath}
\usepackage{xcolor}
% Visible, unmistakable placeholder for any number that must come from the 740-patient grid.
\newcommand{\gridnum}[1]{{\color{red}\textbf{[[GRID: #1]]}}}
% Citation-status marker: \vcite = verified at source (author/year/venue/DOI); \pcite = PENDING.
\newcommand{\vcite}[1]{\cite{#1}}
\newcommand{\pcite}[1]{{\color{orange}\cite{#1}$^{\text{[verify]}}$}}

\begin{document}
\title{Patient-Level vs Nodule-Level Data Partitioning in Pulmonary-Nodule
Classification: A Controlled Paired Study on LIDC-IDRI}
\author{N. Mariotto and co-authors\thanks{Affiliation placeholder.}}
\maketitle

\begin{abstract}
% RESULT-DEPENDENT — structure only; numbers are grid placeholders (D35).
Deep-learning studies on the LIDC-IDRI database routinely report pulmonary-nodule classification
accuracies above 95\% and AUCs above 0.99. Many partition data at the nodule or slice level, so
observations from one patient appear in both training and test folds. We test, \emph{by design},
whether the data-partition unit inflates reported performance: on a single cohort, with identical
labels, preprocessing, architecture, and training, we vary only the partition unit
(patient-grouped vs.\ random) in a controlled paired experiment, repeated over grouped
cross-validation. Across \gridnum{n = repeats$\times$folds} folds and
\gridnum{n. architectures} architectures, random partitioning inflated the leakage-sensitive
metrics by \gridnum{gap log-loss (95\% CI)} in log-loss and \gridnum{gap AUC (95\% CI)} in AUC,
concentrated in \gridnum{calibration vs.\ ranking split}. Patient-level partitioning yields
realistic, substantially lower estimates and should be the methodological standard.
\end{abstract}

% =====================================================================================
\section{Introduction}
% RESULT-INDEPENDENT — safe to write now.
Automated classification of pulmonary nodules on computed tomography (CT) is among the most
studied problems in medical image analysis, and the Lung Image Database Consortium and Image
Database Resource Initiative (LIDC-IDRI) \vcite{armato2011} is its most-used public benchmark.
Reported performance on this database is strikingly high---accuracies above 95\% and areas under
the receiver-operating-characteristic curve (AUC) above 0.99 are common. Such numbers, taken at
face value, would imply near-solved clinical decision support; in practice they have not
translated into deployed tools, which motivates scrutiny of \emph{how} they are obtained.

A recurring methodological choice concerns the unit at which data are partitioned into training
and test sets. LIDC-IDRI contains multiple nodules per patient and multiple annotated slices per
nodule. When partitioning is performed at the nodule or slice level, slices from the same patient
---sharing anatomy, scanner, reconstruction kernel, and acquisition context---can fall on both
sides of the split. A model can then exploit patient-specific context rather than
nodule-intrinsic malignancy signal, a form of data leakage that inflates test performance and
undermines claims of clinical generalisation. The correct unit for a generalisation claim is the
\emph{patient}: every image from a patient must lie entirely within one fold.

Methodological audits in adjacent medical-imaging domains have shown that leakage-free protocols
reduce reported performance by clinically meaningful margins---for example in brain-MRI
prediction \vcite{wen2020} and in COVID-19 chest imaging \vcite{roberts2021}. The pulmonary-nodule
literature, however, has not been subjected to an equivalent \emph{controlled} test: existing
comparisons contrast patient-level results against heterogeneous published numbers that differ in
task, label definition, architecture, and metric convention, so any difference cannot be
attributed to partitioning alone.

\textbf{Gap.} No study on LIDC-IDRI has isolated the partition unit as the single manipulated
variable while holding cohort, labels, preprocessing, architecture, and training fixed, with
repeated cross-validation, uncertainty quantification, and paired significance testing.

\textbf{Contributions.} (i) A controlled paired experiment that isolates the partition unit as the
only changed variable, repeated over grouped cross-validation with per-fold uncertainty and a
paired significance test. (ii) A decomposition of \emph{where} the inflation manifests---ranking
(AUC) vs.\ calibration (log-loss, Brier)---as direct mechanistic evidence of leakage. (iii) A
fully reproducible pipeline: committed partition indices, retained per-sample probabilities,
released evaluation code, and automated integrity gates that make config--code divergence and
cross-experiment mixing impossible to introduce silently.

% =====================================================================================
\section{Related work}
% RESULT-INDEPENDENT — literature; every reference source-verified (author/year/venue/DOI).
\textbf{High-performance nodule classification on LIDC-IDRI.} LIDC-IDRI \vcite{armato2011} is the
standard benchmark for pulmonary-nodule malignancy classification, and reported performance on it is
consistently high. Recent transformer models illustrate the trend: a multitask Swin Transformer
reaches a malignancy ROC-AUC of $0.988$ \vcite{jin2025}, and a compact ViT-Small/16 reaches $92.3\%$
accuracy on a LIDC-IDRI subset \vcite{thakur2026}; convolutional ensembles routinely report
accuracies above $95\%$. We adopt these two transformer families (Swin-Tiny and ViT-Small/16)
precisely because they are the architectures the nodule-classification literature itself uses, and
because together they span hierarchical-windowed and pure-attention designs, letting us test whether
the partitioning effect is specific to convolutional inductive biases.

\textbf{Data partitioning and leakage.} Much of this literature partitions data at the nodule or
slice level. Because LIDC-IDRI contains multiple nodules per patient and multiple annotated slices
per nodule, such partitions place images of one patient on both sides of the train/test boundary.
This is a textbook instance of \emph{data leakage} --- ``the introduction of information about the
\ldots target that should not be legitimately available'' \vcite{kaufman2012} --- and the
appropriate unit for a generalisation claim is therefore the patient, not the nodule or the slice.

\textbf{Leakage audits in medical imaging.} Systematic audits in adjacent domains have quantified the
consequences. In Alzheimer's-disease classification from brain MRI, a reproducible-evaluation study
found that a large fraction of surveyed convolutional studies exhibited evident or unclear data
leakage, and that a leakage-free protocol produced substantially lower, more realistic performance
\vcite{wen2020}. In COVID-19 imaging, a review of models built on chest radiographs and CT concluded
that none was of clinical use, in large part because of data-handling pitfalls including leakage
\vcite{roberts2021}. These audits establish both the phenomenon and the method --- a controlled,
leakage-free re-evaluation; to our knowledge, the present study is the first to apply a controlled
\emph{paired} experiment to the nodule-partitioning question on LIDC-IDRI.

\textbf{Evidence from proper partitioning.} The one comparator that explicitly enforces scan-level
partitioning, the LNDb challenge \vcite{pedrosa2021}, which withholds whole scans for testing,
reports performance in the realistic range rather than the $95$--$99\%$ common under nodule- or
slice-level splits. This is consistent with the hypothesis that the gap is attributable to the
partition unit rather than to architecture, preprocessing, or class definition --- the hypothesis the
present controlled experiment isolates and tests.

% =====================================================================================
\section{Methods}
% RESULT-INDEPENDENT — safe to write now.

\subsection{Dataset, cohort, and patient grouping}
We use LIDC-IDRI \vcite{armato2011} (1{,}018 CT series from 1{,}010 patients) via the \texttt{pylidc}
library \vcite{hancock2016}. Each nodule was read by up to four radiologists; we cluster the
per-radiologist annotations into physical nodules with pylidc's default consensus clustering.

\textbf{Counting (annotation vs.\ nodule).} These two units are reported separately and never
conflated, addressing a specific ambiguity in the prior version. An \emph{annotation} is one
radiologist's reading of a $\geq 3$\,mm nodule (one malignancy score and its contours); a
\emph{nodule} is the physical lesion (a cluster of 1--4 annotations). In LIDC-IDRI this yields
6{,}859 $\geq 3$\,mm annotations that cluster into 2{,}651 physical nodules. Applying the median
label rule (below) leaves 1{,}695 binary-labelled nodules; 956 nodules with a median malignancy of
exactly 3 are excluded as indeterminate. The labelling and grouping unit is always the nodule.

\textbf{Cohort.} The principal cohort keeps every labelled nodule with at least one annotator
(comparable to the criticised literature), giving 11{,}026 slice-samples (5{,}504 malignant /
5{,}522 benign) across 740 patients at the slice level, and 1{,}695 nodules (545 malignant /
1{,}150 benign) at the nodule level. A high-agreement cohort ($\geq 3$ annotators; 6{,}973 slices /
879 nodules) is reported as a label-noise sensitivity analysis. Slices per patient range from 1 to
97 (median 11): this within-patient correlation is precisely why the partition unit matters and
motivates patient-level grouping.

\textbf{Patient grouping.} All partitioning is at the level of the \emph{patient}: folds are formed
by grouping on patient identity, and we assert \emph{and export} that no patient identifier appears
in more than one of the train/validation/test partitions of any fold. The assertion runs at split
time and is re-verified on the committed indices (Section~\ref{sec:repro}); we measured zero
patient overlap across train/test/validation in every patient-arm fold.

\subsection{Label definition}
Malignancy is binarised from the radiologists' 1--5 scores using the median across annotators:
scores 4--5 are malignant, 1--2 benign, and score-3 (indeterminate) nodules are excluded. This
follows established practice and keeps the label identical across the two partitioning arms so it
cannot confound the contrast.

\subsection{Preprocessing}
For each nodule we take the CT volume in Hounsfield units, clip to $[-1000, 400]$ HU, and extract
a region of interest from the 50\% consensus mask ($\geq 2$ of 4 annotators)
\vcite{kubota2011}. The ROI is squared, padded, and resized to $256\times256$ and replicated to
three channels for ImageNet-pretrained backbones. Preprocessing is \emph{identical} across the two
arms; image enhancement (e.g.\ CLAHE) is analysed separately as a controlled factor rather than
folded into the main contrast.

\subsection{The controlled paired design}
The experiment is a matched pair. In arm~A (\emph{patient}) folds are formed by grouped
cross-validation on patient identity; in arm~B (\emph{random}) folds are formed by shuffling the
same samples ignoring patient identity. \textbf{Only the partition unit differs}---cohort, labels,
preprocessing, architecture, optimiser, schedule, and evaluation set size are held fixed. The
design is repeated over $R$ repeats $\times$ $K$ folds (grouped), giving
$n = R\times K \ge 10$ matched pairs, and is run on both the slice and the nodule granularity
(a $2\times2$ of partition-unit $\times$ sample-unit). The leakage-advantage gap is reported per
(architecture, sample-unit).

A central property of this matched design is that any factor shared identically by both arms
cancels in the arm-to-arm difference. In particular, the ImageNet-pretrained initialisation, the
preprocessing, and the architecture are common to both arms; they may raise or lower the absolute
performance of both, but they cannot create a difference between them and therefore cannot bias the
reported gap. This is a substantive advantage over a single-number study, whose absolute metric
would inherit any such shared bias --- including any hidden overlap between a pretrained model's
(often undisclosed) training data and the evaluation set. Our claims are consequently stated as the
gap, not as absolute performance.

\subsection{Formalising the leakage}
Let the dataset be $\mathcal{D}=\{(x_i,y_i,p_i,n_i)\}_{i=1}^{N}$, where $x_i$ is a slice, $y_i\in\{0,1\}$
its malignancy label, $p_i$ its patient and $n_i$ its nodule. A partition is a map
$\pi:\{1,\dots,N\}\to\{\mathrm{train},\mathrm{val},\mathrm{test}\}$. The two arms differ \emph{only}
in $\pi$. The patient-grouped arm~A uses a $\pi_A$ that is constant on patients,
$p_i=p_j \Rightarrow \pi_A(i)=\pi_A(j)$, so the test set is patient-disjoint from train:
$\{p_i:\pi_A(i)=\mathrm{test}\}\cap\{p_i:\pi_A(i)=\mathrm{train}\}=\varnothing$. The random arm~B
assigns samples independently of $p_i$, so patients may straddle the boundary.

Define the patient-leakage rate of a partition as
\begin{equation}
L(\pi)=\frac{\bigl|\{i:\pi(i)=\mathrm{test},\ p_i\in\mathcal{P}_{\mathrm{train}}(\pi)\}\bigr|}
{\bigl|\{i:\pi(i)=\mathrm{test}\}\bigr|},\qquad
\mathcal{P}_{\mathrm{train}}(\pi)=\{p_j:\pi(j)=\mathrm{train}\}.
\end{equation}
By construction $L(\pi_A)=0$; empirically $L(\pi_B)\approx0.99$. For a metric $M$, the
leakage-advantage gap is $\Delta_M=M(\pi_B)-M(\pi_A)$ (sign reversed for log-loss/Brier). Because
the arms share cohort, labels, preprocessing, architecture and training, any contribution $c$
common to both cancels: if $M(\pi)=g(\pi)+c$ then $\Delta_M=g(\pi_B)-g(\pi_A)$, independent of $c$;
thus $\Delta_M$ isolates the partition unit.

Slice-level leakage decomposes into two disjoint routes,
\begin{equation}
\{i\in\mathrm{test}:p_i\in\mathcal{P}_{\mathrm{train}}\}=
\underbrace{\{i:\exists\,j\in\mathrm{train},\,n_i=n_j\}}_{\text{within-nodule redundancy}}\ \cup\
\underbrace{\{i:\exists\,j\in\mathrm{train},\,p_i=p_j,\,n_i\neq n_j\}}_{\text{cross-nodule, same patient}},
\end{equation}
and aggregating to the nodule level (one sample per nodule) eliminates the first term, so the
nodule-level gap $\Delta_M^{\mathrm{nod}}$ reflects only the patient-level route --- the component a
3D architecture would also incur (cf.\ Reviewer~4's 2D-vs-3D question).

\subsection{Architectures and training}
We evaluate \gridnum{n. architectures} ImageNet-pretrained backbones spanning convolutional
families and a transformer/hybrid baseline (satisfying the reviewers' request for
transformer/SOTA models). Each run trains for at most 30 epochs with early stopping (patience 10)
on validation loss, selecting the best-validation-loss checkpoint; a class-weighted cross-entropy
loss accommodates imbalance. Training uses mixed precision and gradient accumulation to fit an
8\,GB GPU. The exact configuration is frozen and its hash is stamped into every run so that no run
trained under a different configuration can be folded into an average (Section~\ref{sec:repro}).

\subsection{Evaluation and statistics}
Per fold we compute AUC, log-loss, and Brier score (and accuracy, precision, recall, F1 at the
0.5 threshold), at the nodule level (mean of a nodule's slice probabilities) and at the slice
level. Metrics are computed \emph{within} each fold and averaged across folds; predictions are
never pooled across folds. Because cross-validation fold estimates share training data, confidence
intervals use the Nadeau--Bengio correction \vcite{nadeau2003} rather than the naive standard
error. The paired A-vs-B gap is tested with the Wilcoxon signed-rank test; we report its exact
attainable minimum $p$-value at the given $n$ so that a floor value is not mistaken for a
near-miss. Per-metric bootstrap confidence intervals resample the independent unit (the patient).
Within-arm architecture comparisons (identical test set) use McNemar's test; it is \emph{not}
applied across arms, whose test sets differ.

\subsection{Relationship to the prior submission}
An earlier version of this work compared three heterogeneous scenarios (single-site multiclass;
binary malignancy; an end-to-end segmentation--classification pipeline) and attributed a monotonic
performance decline to increasing realism. As Reviewer~3 correctly observed, ``the monotonic
decline \ldots is \emph{not} evidence of increasing realism because the three experiments change
the dataset, classification target, processing pipeline, number of models, and evaluation protocol
simultaneously,'' and ``a valid analysis would require controlled paired experiments in which only
the partitioning unit changes, repeated across several random seeds or grouped cross-validation
folds, with uncertainty estimates and statistical comparisons.'' We agree. We have therefore
\emph{removed} the three-scenario comparison and replaced it with the controlled paired experiment
described above, which isolates the partition unit as the single manipulated variable. This
directly addresses the central objection rather than patching the confounded comparison.

\subsection{Reproducibility and integrity safeguards}\label{sec:repro}
% RESULT-INDEPENDENT — describes machinery, not outcomes.
Every reported quantity is produced by released code from artifacts on disk; no value is
hand-transcribed. Partition indices are committed; per-sample probabilities are retained;
train/validation history is stored per run. A hard coverage gate refuses to train unless every
split sample has a preprocessed image. Writes are atomic. A content-validated skip check counts a
run as ``done'' only if its weights load, its probabilities are finite and length-matched to the
test split, and its stamped configuration hash matches the frozen configuration. A config--code
contract asserts that every configuration key is read by code (or explicitly declared inert), so
the manuscript cannot describe behaviour the code does not implement. A final gate asserts that the
entire grid shares one configuration, epoch budget, and patience, so results are never averaged
across mixed experiments.

The reported results are reproducible from the released artifacts: the committed partition indices,
the retained per-sample probabilities, and the evaluation code regenerate every table and figure
without retraining. Training itself uses standard non-deterministic GPU acceleration (cuDNN
autotuning), so exact bit-wise retraining is not guaranteed; this is by design, as the reported
statistics are computed over \gridnum{repeats} repeats $\times$ \gridnum{folds} grouped folds with
uncertainty estimates, and the run-to-run stochasticity of training is thereby captured in the
reported confidence intervals rather than suppressed. Reproducibility of the reported numbers
(from artifacts) is thus decoupled from bit-wise reproducibility of the training process.

% =====================================================================================
\section{Results}
% RESULT-DEPENDENT — DO NOT WRITE until the 740 grid completes and verify_grid_consistency passes
% (D35). Structure only; all numbers are placeholders sourced from outputs/results/*.
\gridnum{Table: per-architecture leakage gap in AUC / log-loss / Brier, nodule and slice level,
with Nadeau--Bengio 95\% CIs and Wilcoxon p (with attainable floor)}

\gridnum{Figure: learning curves, patient vs.\ random, mean$\pm$sd over folds --- the leakage
mechanism (random-arm validation keeps improving = memorisation; patient-arm converges early)}

\gridnum{Figure/Table: confusion matrices, patient vs.\ random arm, per architecture}

\gridnum{Per-model metric tables (I3); within-arm McNemar (B5)}

% =====================================================================================
\section{Discussion}
% STRUCTURE + result-independent parts only. Mechanism numbers are placeholders (D35).
\textbf{Principal finding.} \gridnum{one-sentence result with gap$\pm$CI and the honest hedging}.

\textbf{Mechanism.} The inflation is concentrated in calibration (log-loss, Brier) more than in
ranking (AUC): \gridnum{quantified calibration-vs-ranking split, re-derived from the 740 grid}.
This is consistent with leakage acting through patient-specific context that sharpens confidence
on seen patients rather than improving nodule discrimination. The learning curves make the
mechanism visible: under random partitioning the validation loss tracks training loss and keeps
falling (the leaked validation set is memorised alongside training), whereas under patient
partitioning validation loss diverges upward from training loss, the expected honest
generalisation gap.

\textbf{The gap is not a test-set artefact.} Because the two arms have different test sets by
construction, a natural objection is that the gap reflects test-set composition rather than
leakage. We control for this by re-scoring both arms on the \emph{identical} rows (the same nodules
and slices present in both arms' test sets), holding composition, prevalence, size, and aggregation
depth fixed: the gap \gridnum{persists / grows under the matched control}. Composition therefore
does not explain the effect.

\textbf{Why the partition unit, and not stratification, governs the variance.} In grouped
cross-validation the fold-to-fold variance is driven by \emph{which patients} (hence which cases)
fall in a fold, not by the class proportion. We show this directly: the intra-class correlation of
the label within a patient is \gridnum{ICC}, so patients behave as partial label-blocks; an
off-the-shelf stratified grouped split does not reduce fold-prevalence spread here
(\gridnum{stratified vs.\ grouped spread}), and even a near-perfect balancer changes the
fold-to-fold metric variance by only \gridnum{small \%}. The practical consequence is that
\emph{only $n$} (repeats $\times$ folds), not stratification, tightens the interval --- which is
why the design is repeated rather than stratified, and why single-split studies cannot quantify the
gap at all.

\textbf{Methodological implication.} Reporting nodule- or slice-level performance without
patient-level partitioning does not measure clinical generalisation; it measures, in part, a
model's ability to recognise patients it has already seen. The gap we report is a floor on that
inflation under a single, clean manipulation. We therefore recommend patient-level partitioning as
the default for pulmonary-nodule classification, and that studies state the partition unit
explicitly and release executable partition indices so the choice is auditable.

\textbf{Limitations.} The study uses 2D backbones on volumetric CT, mirroring the common practice
it audits; a 3D interaction is future work. Labels are radiologist-consensus rather than
histopathology; this label noise is identical across arms and so does not bias the paired
contrast. The analysis is confined to LIDC-IDRI; external, multi-institutional validation is
valuable but out of scope for a methodological audit and is stated as future work. We validated the
pipeline on a 250-patient subset before scaling; those subset numbers are a pipeline check, not a
study result, and are not reported as such (D35).

\textbf{Relation to prior audits.} The magnitude of the effect is consistent with leakage gaps
documented in adjacent medical-imaging domains \vcite{wen2020,roberts2021}.

% =====================================================================================
\section{Conclusion}
Partitioning pulmonary-nodule data at the nodule or slice level allows patient information to leak
across the train/test boundary and inflates reported performance. In a controlled paired
experiment that changes only the partition unit, patient-level evaluation produces realistic and
substantially lower estimates. Patient-level partitioning should be the default methodological
standard, and studies should report the partition unit explicitly.

% =====================================================================================
% CREdiT / author contributions (IEEE Access checklist) — PLACEHOLDER, researcher to complete.
% \section*{Author Contributions} Using CRediT roles per author (Conceptualization, Methodology,
% Software, Validation, Formal analysis, Investigation, Data curation, Writing - original draft,
% Writing - review \& editing, Visualization, Supervision). REQUIRED by the checklist; fill with the
% real per-author split.

\section*{Acknowledgment}
% AI-USE DISCLOSURE — required by the IEEE policy "Author Guidelines for Artificial Intelligence
% (AI)-Generated Text" (open.ieee.org): AI-generated content (text, figures, images, AND code) must
% be disclosed in the Acknowledgments, identifying the system, the sections, and the level of use.
% Written to reflect the ACTUAL, larger role of AI in this rebuild (not the original's narrow note).
The authors used a large language model (Claude, Anthropic), operated through the agentic coding
assistant Claude Code, during the preparation of this work. Under the authors' direction, the tool
assisted in: (i) implementing the experimental-pipeline code (metadata construction, preprocessing,
partition-index generation, model training, and evaluation); (ii) implementing the
statistical-analysis and figure-generation code; (iii) drafting portions of the manuscript text,
including the Introduction, Related Work, Methods, and Discussion, and the response to reviewers;
and (iv) assisting in verifying the bibliographic details of the cited references against primary
sources. The authors specified the experimental design and every methodological decision, reviewed
and verified all AI-assisted code and text, and independently confirmed each reported result against
the artifacts produced by the released code. The scientific content, methodological decisions, data,
results, interpretations, and conclusions are the authors' own; the authors take full responsibility
for the integrity and accuracy of the manuscript. No AI system is an author of this work.

% =====================================================================================
% CITATION VERIFICATION STATUS (D-doctrine + IEEE checklist "accurate bibliographic details":
% verify author/year/venue/DOI at source before a ref is treated as verified).
% ALL CITATIONS VERIFIED AT SOURCE 2026-07-20 — no \pcite remains.
%     armato2011  — Armato et al. LIDC-IDRI. Med. Phys. 38(2):915-931, 2011. DOI 10.1118/1.3528204.
%     hancock2016 — Matthew C. Hancock, Jerry F. Magnan. Lung nodule malignancy classification using
%                   only radiologist-quantified image features ... probing the LIDC dataset. J. Med.
%                   Imaging 3(4):044504, 2016. DOI 10.1117/1.jmi.3.4.044504. [Crossref, full author list]
%     kubota2011  — consensus-mask basis (D19). Med. Image Anal. 15(1):133-154.
%                   DOI 10.1016/j.media.2010.08.005.
%     nadeau2003  — Nadeau & Bengio, Inference for the Generalization Error. Machine Learning
%                   52:239-281, 2003. DOI 10.1023/A:1024068626366.
%     wen2020     — Wen, Thibeau-Sutre, Diaz-Melo, Samper-Gonzalez, Routier, Bottani, Dormont,
%                   Durrleman, Burgos, Colliot. CNNs for classification of Alzheimer's disease:
%                   overview and reproducible evaluation. Med. Image Anal. 63:101694, 2020.
%                   DOI 10.1016/j.media.2020.101694. [Crossref]
%     roberts2021 — Roberts et al., Common pitfalls and recommendations ... COVID-19 ... chest
%                   radiographs and CT. Nat. Mach. Intell. 3:199-217, 2021. DOI 10.1038/s42256-021-00307-0.
%     pedrosa2021 — Pedrosa, Aresta, ..., Campilho. LNDb challenge on automatic lung cancer patient
%                   management. Med. Image Anal. 70:102027, 2021. DOI 10.1016/j.media.2021.102027. [Crossref]
%     setio2017   — LUNA16 (criterion only, if E2E ever revived). Med. Image Anal. 42:1-13.
%                   DOI 10.1016/j.media.2017.06.015; PMID 28732268.
%     kaufman2012 — Kaufman, Rosset, Perlich, Stitelman. Leakage in Data Mining: Formulation,
%                   Detection, and Avoidance. ACM TKDD 6(4):1-21, 2012. DOI 10.1145/2382577.2382579. [Crossref]
%     jin2025     — Jin H, Yu C, Zhang J, Zheng R, Fu Y, Zhao Y. Multitask Swin Transformer for
%                   classification and characterization of pulmonary nodules in CT images. Quant.
%                   Imaging Med. Surg. 15(3):1845-1861, 2025. DOI 10.21037/qims-24-1619. [publisher page]
%                   (cite the PUBLISHED version, NOT the SSRN preprint 10.2139/ssrn.4597429)
%     thakur2026  — Thakur, Chouksey, Sharma, Chopra, Sadotra, Kumar. Binary classification of lung
%                   cancer using vision transformer models on CT images. Discover Computing 29(1):66,
%                   2026. DOI 10.1007/s10791-026-09931-z. [Crossref]
% ALL 11 REFERENCES NOW FULLY VERIFIED AT SOURCE (Crossref / publisher page): author list,
% year, venue, and DOI confirmed for every one. Reference gate CLOSED (2026-07-21).
% When compiling the .bib: list references in ORDER OF CITATION (IEEE checklist).
% Cautionary note: jin2025 (MTST Swin) — the SSRN preprint (10.2139/ssrn.4597429) lists 5 authors,
% the PUBLISHED QIMS version (10.21037/qims-24-1619) lists 6 (adds Zhang J). Citing the preprint
% would have produced a WRONG author list — the exact defect class of the 3 hallucinated refs in the
% original submission. Always cite the published, source-verified version.
\end{document}
